\section{Evaluation Metrics}
\label{sec:metrics}

The evaluation compares a predicted adjacency matrix \(\hat{A}_{m,s}\) with
the ground-truth adjacency matrix \(A^\ast\). Both matrices are binary and use
the same six-feature ordering. A value of 1 at row \(i\), column \(j\) means
that the source feature \(i\) has a directed relationship to the target
feature \(j\). The supervisor notebook binarizes matrices with a threshold of
0.5, which is equivalent to reading the saved 0/1 adjacency matrices
directly.

The metric suite contains eight values. Four metrics are edge-classification
metrics over the flattened adjacency matrix. Four metrics summarize structural
or matrix-level agreement. Table~\ref{tab:graph-metrics} gives the evaluation
role of each metric.

\begin{table}[!htbp]
  \centering
  \small
  \setlength{\tabcolsep}{3pt}
  \begin{tabularx}{\textwidth}{@{}L{36mm}L{26mm}Y@{}}
    \toprule
    Metric & Direction & Interpretation \\
    \midrule
    Precision & Higher is better & Fraction of predicted edges that are present in the ground truth. \\
    Recall & Higher is better & Fraction of ground-truth edges recovered by the model output. \\
    F1 Score & Higher is better & Harmonic mean of precision and recall; the headline edge-recovery balance. \\
    SHD & Lower is better & Structural Hamming distance; count of entries where predicted and ground-truth adjacency matrices differ. \\
    FDR & Lower is better & False discovery rate; fraction of predicted edges that are false positives. \\
    Pearson Correlation & Higher is better & Correlation between flattened predicted and ground-truth adjacency matrices, with undefined constant cases mapped to 0. \\
    Hamming Distance & Lower is better & Fraction of flattened adjacency entries that differ. \\
    Frobenius Norm & Lower is better & Matrix norm of the binary difference \(P-G\); numerically equal to \(\sqrt{\mathrm{SHD}}\) for binary matrices. \\
    \bottomrule
  \end{tabularx}
  \caption{Graph-recovery metrics used in the supervisor-provided evaluation matrix. The suite combines edge-level classification metrics with matrix-level structural agreement measures.}
  \label{tab:graph-metrics}
\end{table}

\subsection{Confusion Counts}

Let \(\hat{A}_{ij}\) be the predicted adjacency entry for one model-strategy
cell and \(A^\ast_{ij}\) the ground-truth entry. The flattened matrices define
four counts:

\begin{align}
  TP &= \sum_{i,j} \mathbb{1}[\hat{A}_{ij}=1 \wedge A^\ast_{ij}=1],\\
  FP &= \sum_{i,j} \mathbb{1}[\hat{A}_{ij}=1 \wedge A^\ast_{ij}=0],\\
  FN &= \sum_{i,j} \mathbb{1}[\hat{A}_{ij}=0 \wedge A^\ast_{ij}=1],\\
  TN &= \sum_{i,j} \mathbb{1}[\hat{A}_{ij}=0 \wedge A^\ast_{ij}=0].
\end{align}

These counts are computed over all matrix cells, including non-edge cells.
Because the graph is sparse, metrics must be interpreted carefully. Accuracy
is not reported as a headline metric because a model can obtain many true
negatives simply by predicting few or no edges. Precision, recall, F1, and
structural error are more informative for this task.

\subsection{Precision, Recall, and F1 Score}

Precision measures how many predicted edges are correct:

\begin{equation}
  \mathrm{Precision} = \frac{TP}{TP + FP}.
\end{equation}

Recall measures how many ground-truth edges are recovered:

\begin{equation}
  \mathrm{Recall} = \frac{TP}{TP + FN}.
\end{equation}

The F1 score balances precision and recall:

\begin{equation}
  \mathrm{F1} =
  \frac{2 \cdot \mathrm{Precision} \cdot \mathrm{Recall}}
       {\mathrm{Precision} + \mathrm{Recall}}.
\end{equation}

The notebook uses zero-division handling: if a denominator is zero, the
corresponding score is set to 0. This matters for cells where a model produces
no correct predicted edges or no predicted edges after parsing.

\subsection{Structural Hamming Distance}

Structural Hamming distance is the number of adjacency entries that differ:

\begin{equation}
  \mathrm{SHD} = \sum_{i,j} \mathbb{1}[\hat{A}_{ij} \neq A^\ast_{ij}].
\end{equation}

In this binary directed-graph setting, SHD counts both missed ground-truth
edges and extra predicted edges. It is useful because it reports absolute
structural error rather than only a precision-recall balance. A low SHD means
the predicted matrix is close to the target matrix; a high SHD means many
edge/non-edge decisions differ.

\subsection{False Discovery Rate}

False discovery rate is the fraction of predicted edges that are not in the
ground truth:

\begin{equation}
  \mathrm{FDR} = \frac{FP}{TP + FP}.
\end{equation}

When no edges are predicted, the supervisor notebook sets FDR to 0. This
convention avoids division by zero but requires careful interpretation:
zero FDR does not necessarily indicate a useful model if recall and F1 are
also zero.

\subsection{Pearson Correlation}

Pearson correlation is computed between the flattened binary predicted matrix
and the flattened binary ground-truth matrix. It captures linear agreement
between edge patterns. The notebook maps undefined cases to 0 when either
flattened vector is constant or the correlation is not a number. This
conservative fallback prevents degenerate outputs from producing undefined
metric values.

\subsection{Hamming Distance and Frobenius Norm}

Hamming distance is the fraction of flattened matrix entries that differ:

\begin{equation}
  \mathrm{Hamming}(\hat{A},A^\ast) =
  \frac{1}{n^2}\sum_{i,j}\mathbb{1}[\hat{A}_{ij} \neq A^\ast_{ij}].
\end{equation}

For the six-feature task, \(n=6\), so the denominator is 36. Hamming distance
is therefore a normalized version of structural error.

The Frobenius norm is:

\begin{equation}
  \lVert \hat{A}-A^\ast\rVert_F =
  \sqrt{\sum_{i,j}(\hat{A}_{ij}-A^\ast_{ij})^2}.
\end{equation}

Because the matrices are binary, this value is the square root of the number
of differing entries. It provides a matrix-error magnitude that is consistent
with SHD while using a standard matrix norm.

\section{Metric Interpretation}

No single metric is sufficient. Precision can be high when a model predicts
few edges, recall can be high when a model predicts many edges, and F1 can
hide whether the errors are mainly false positives or false negatives. SHD,
Hamming distance, and Frobenius norm expose total matrix error. FDR exposes
how much of the predicted graph is spurious. Pearson correlation provides an
additional matrix-level agreement signal.

The results chapter therefore reports rankings and interpretation across the
metric suite rather than treating one number as a universal quality score. F1
is useful as the primary edge-recovery balance, but the thesis also examines
false discovery and structural error because fleet-monitoring graphs can be
misleading when they contain plausible but unsupported directed links.
